\documentclass{article}
\usepackage{amsmath, amssymb, amsthm}
\usepackage{hyperref}
\usepackage{geometry}
\geometry{margin=1in}

\title{Erd\H{o}s Problem \#414}
\date{Last updated: 2025-08-31}
\author{Source: erdosproblems.com}

\begin{document}
\maketitle

\noindent\textbf{Status:} OPEN \quad \textbf{No prize}

\noindent\textbf{Tags:} number theory, iterated functions

\medskip
\noindent\textbf{Problem Statement:}

Let $h_1(n)=h(n)=n+\tau(n)$ (where $\tau(n)$ counts the number of divisors of $n$) and $h_k(n)=h(h_{k-1}(n))$. Is it true, for any $m,n$, there exist $i$ and $j$ such that $h_i(m)=h_j(n)$?

\bigskip
\noindent\textbf{Remarks:}

Asked by Spiro. That is, there is (eventually) only one possible sequence that the iterations of $n\mapsto h(n)$ can settle on. Erd\H{o}s and Graham believed the answer is yes. Similar questions can be asked by the iterates of many other functions. See also [412] and [413].

\bigskip
\noindent\small{Source: \url{https://www.erdosproblems.com/414}}
\end{document}
